\documentclass[11pt,a4paper]{article}
\usepackage[utf8]{inputenc}
\usepackage[T1]{fontenc}
\usepackage[catalan]{babel}
\usepackage{lmodern}
\usepackage[sfdefault,lining]{FiraSans}
\usepackage[a4paper,top=2.3cm,bottom=2.2cm,left=1.8cm,right=1.8cm,headheight=1.3cm,headsep=0.4cm,footskip=1cm]{geometry}
\usepackage{amsmath,amssymb}
\usepackage{icomma}
\usepackage{graphicx}
\usepackage[table]{xcolor}
\usepackage{tikz}
\usepackage{fontawesome5}
\usepackage{ccicons}
\usepackage{fancyhdr}
\usepackage{lastpage}
\usepackage{tasks}
\usepackage{needspace}
\usepackage[most]{tcolorbox}
\usepackage{array}
\usepackage[hidelinks,pdftitle={Fitxa d'exercicis de la unitat 10: Proporcionalitat. Teorema de Tales. 3r d'ESO},pdfsubject={Matemàtiques 3r ESO},pdfkeywords={raó, proporció, segments proporcionals, teorema de Tales, posició de Tales, semblança, raó de semblança, criteris de semblança, escala}]{hyperref}

\definecolor{navy}{HTML}{1B2748}
\definecolor{accent}{HTML}{E8603A}
\definecolor{teal}{HTML}{1F9E90}
\definecolor{soft}{HTML}{EEF1F8}
\definecolor{muted}{HTML}{5D6685}

\setlength{\parindent}{0pt}
\setlength{\parskip}{3pt}
\settasks{label=\alph*), label-format=\bfseries\color{navy}, label-width=1.3em, item-indent=1.8em, column-sep=0.6em, after-item-skip=5pt, before-skip=3pt, after-skip=4pt}

% ---------- capçalera i peu ----------
\newcommand{\llicencia}{\href{https://creativecommons.org/licenses/by-nc-sa/4.0/deed.ca}{\ccbyncsa}\ \ Llicència Creative Commons Reconeixement-NoComercial-CompartirIgual 4.0 Internacional (CC BY-NC-SA 4.0)}
\pagestyle{fancy}
\fancyhf{}
\fancyhead[L]{\raisebox{-0.35\height}{\includegraphics[height=0.95cm]{logo-gva.png}}}
\fancyhead[C]{\small\color{muted}Fitxa d'exercicis · Unitat 10: Proporcionalitat. Teorema de Tales · 3r d'ESO}
\fancyhead[R]{\raisebox{-0.35\height}{\includegraphics[height=0.95cm]{logo-vermellar.png}}}
\fancyfoot[L]{\scriptsize\color{muted}\llicencia}
\fancyfoot[R]{\small\color{muted}\thepage}
\renewcommand{\headrulewidth}{0.4pt}
\renewcommand{\headrule}{\color{muted!40}\hrule height \headrulewidth}
\fancypagestyle{primera}{\fancyhf{}\renewcommand{\headrulewidth}{0pt}
  \fancyfoot[L]{\scriptsize\color{muted}\llicencia}
  \fancyfoot[R]{\small\color{muted}\thepage}}

% ---------- nivells i exercicis ----------
\newcommand{\estrelles}[1]{{\color{accent}\ifcase#1\or\faStar\faStar[regular]\faStar[regular]\or\faStar\faStar\faStar[regular]\or\faStar\faStar\faStar\fi}}
\newcounter{exer}
\newcommand{\ex}[1]{\par\needspace{4\baselineskip}\medskip\stepcounter{exer}\noindent{\color{navy}\bfseries\theexer.}\ \ {\small\estrelles{#1}}\ \ }

\newcommand{\seccio}[2]{\par\needspace{10\baselineskip}\bigskip
  \begin{tcolorbox}[colback=navy,colframe=navy,coltext=white,arc=2mm,boxrule=0pt,left=3mm,right=3mm,top=1.2mm,bottom=1.2mm]
  {\large\bfseries #1\quad #2}
  \end{tcolorbox}\vspace{-2pt}}

\newcolumntype{C}[1]{>{\centering\arraybackslash}p{#1}}
\usepackage{pifont}

% quadre de recordatori (només a la fitxa d'exercicis)
\newcommand{\recorda}[1]{\par\smallskip\begin{tcolorbox}[colback=accent!7,colframe=accent!55,arc=2mm,boxrule=0.5pt,left=2.5mm,right=2.5mm,top=1mm,bottom=1mm]\small{\color{accent}\faLightbulb[regular]}\ \textbf{Recorda:}\ #1\end{tcolorbox}}
\newcommand{\res}[1]{#1}
\newcommand{\rest}[1]{#1}

\def\fmneed{14}
% =====================================================================
% Capçalera FentMates de les fitxes: un color per a cada curs.
% Els colors són els de la pàgina inicial (LEVELS a assets/js/materials.js);
% el logo ve de _eines/logo/pdf/ (python3 _eines/logo/genera.py).
%
% Ús, al final del preàmbul de la unitat:
%   % =====================================================================
% Capçalera FentMates de les fitxes: un color per a cada curs.
% Els colors són els de la pàgina inicial (LEVELS a assets/js/materials.js);
% el logo ve de _eines/logo/pdf/ (python3 _eines/logo/genera.py).
%
% Ús, al final del preàmbul de la unitat:
%   % =====================================================================
% Capçalera FentMates de les fitxes: un color per a cada curs.
% Els colors són els de la pàgina inicial (LEVELS a assets/js/materials.js);
% el logo ve de _eines/logo/pdf/ (python3 _eines/logo/genera.py).
%
% Ús, al final del preàmbul de la unitat:
%   \input{../comu/fentmates.tex}
%   \fmcurs{1eso}
% i, al començament del document:
%   \capcalera{Matemàtiques}{1r ESO}{Fitxa d'exercicis · Unitat 2}{Divisibilitat}
% Opcional abans de \input: \def\fmneed{14} (espai mínim abans de cada secció).
% =====================================================================
\makeatletter
\def\fm@logodir{../../logo/pdf}
\providecommand{\fmneed}{10}
\def\fm@cursnom{}

% clau del curs, color i color del text damunt del color
\newcommand{\fm@curs}[3]{\@namedef{fm@color@#1}{#2}\@namedef{fm@text@#1}{#3}}
\fm@curs{1eso}{E6B84A}{navy}
\fm@curs{2eso}{3B82F6}{white}
\fm@curs{3eso}{FF7A50}{white}
\fm@curs{4eso-a}{A78BFA}{white}
\fm@curs{4eso-b}{6D28D9}{white}
\fm@curs{1batx-mi}{9B1C1C}{white}
\fm@curs{1batx-ccss}{EF6B6B}{white}
\fm@curs{2batx-mii}{1E6B47}{white}
\fm@curs{2batx-ccss}{2BB3A3}{white}

\newcommand{\fmcurs}[1]{%
  \@ifundefined{fm@color@#1}{\PackageError{fentmates}{Curs desconegut: #1}{}}{}%
  \edef\fm@tmp{\noexpand\definecolor{curs}{HTML}{\@nameuse{fm@color@#1}}%
               \noexpand\colorlet{curstext}{\@nameuse{fm@text@#1}}}\fm@tmp
  \colorlet{cursfosc}{curs!72!black}%
  \def\fm@logo{\fm@logodir/vertical-#1.pdf}%
  \def\fm@marca{\fm@logodir/marca-#1.pdf}}

% ---------- capçalera de la primera pàgina ----------
% #1 matèria (petit), #2 curs, #3 tipus de document i unitat, #4 títol
\newcommand{\capcalera}[4]{%
  \gdef\fm@cursnom{#2}%
  \par\noindent
  \begin{tcolorbox}[colback=curs,colframe=curs,boxrule=0pt,arc=2.2mm,width=3.5cm,height=2.15cm,
      valign=center,left=3mm,right=2mm,top=1mm,bottom=1mm,nobeforeafter,box align=center]
    \color{curstext}{\footnotesize #1}\par\vspace{1pt}{\fontsize{19}{21}\selectfont\bfseries #2\par}%
  \end{tcolorbox}\hspace{5mm}%
  \begin{minipage}[c]{\dimexpr\linewidth-3.5cm-5mm-2.9cm\relax}
    {\color{cursfosc}\small\bfseries\MakeUppercase{#3}}\par\vspace{2pt}
    {\color{navy}\fontsize{20}{23}\selectfont\bfseries #4\par}
  \end{minipage}\hfill
  \begin{minipage}[c]{2.6cm}\raggedleft\includegraphics[width=2.6cm]{\fm@logo}\end{minipage}\par
  \vspace{2pt}}

% ---------- barres de secció del color del curs ----------
\renewcommand{\seccio}[2]{\par\needspace{\fmneed\baselineskip}\bigskip
  \begin{tcolorbox}[colback=curs,colframe=curs,coltext=curstext,arc=2mm,boxrule=0pt,left=3mm,right=3mm,top=1.2mm,bottom=1.2mm]
  {\large\bfseries #1\quad #2}
  \end{tcolorbox}\vspace{-2pt}}

% ---------- capçalera i peu de les pàgines ----------
% primera pàgina: logos institucionals al peu
\fancypagestyle{primera}{\fancyhf{}\renewcommand{\headrulewidth}{0pt}
  \fancyfoot[L]{\raisebox{-0.32\height}{\includegraphics[height=0.78cm]{logo-gva.png}}\hspace{5mm}%
    \raisebox{-0.32\height}{\includegraphics[height=0.78cm]{logo-vermellar.png}}}
  \fancyfoot[R]{\scriptsize\color{muted}\href{https://creativecommons.org/licenses/by-nc-sa/4.0/deed.ca}{\ccbyncsa}\ \ CC BY-NC-SA 4.0 · fentmates.org\qquad\small\thepage}}
% la resta: el símbol a l'esquerra i el curs a la dreta, amb el títol al centre com abans
\fancyhead[L]{\raisebox{-0.3\height}{\includegraphics[height=0.62cm]{\fm@marca}}}
\fancyhead[R]{\small\bfseries\color{cursfosc}\fm@cursnom}
\renewcommand{\headrule}{\color{curs!45}\hrule height \headrulewidth}
\makeatother

%   \fmcurs{1eso}
% i, al començament del document:
%   \capcalera{Matemàtiques}{1r ESO}{Fitxa d'exercicis · Unitat 2}{Divisibilitat}
% Opcional abans de \input: \def\fmneed{14} (espai mínim abans de cada secció).
% =====================================================================
\makeatletter
\def\fm@logodir{../../logo/pdf}
\providecommand{\fmneed}{10}
\def\fm@cursnom{}

% clau del curs, color i color del text damunt del color
\newcommand{\fm@curs}[3]{\@namedef{fm@color@#1}{#2}\@namedef{fm@text@#1}{#3}}
\fm@curs{1eso}{E6B84A}{navy}
\fm@curs{2eso}{3B82F6}{white}
\fm@curs{3eso}{FF7A50}{white}
\fm@curs{4eso-a}{A78BFA}{white}
\fm@curs{4eso-b}{6D28D9}{white}
\fm@curs{1batx-mi}{9B1C1C}{white}
\fm@curs{1batx-ccss}{EF6B6B}{white}
\fm@curs{2batx-mii}{1E6B47}{white}
\fm@curs{2batx-ccss}{2BB3A3}{white}

\newcommand{\fmcurs}[1]{%
  \@ifundefined{fm@color@#1}{\PackageError{fentmates}{Curs desconegut: #1}{}}{}%
  \edef\fm@tmp{\noexpand\definecolor{curs}{HTML}{\@nameuse{fm@color@#1}}%
               \noexpand\colorlet{curstext}{\@nameuse{fm@text@#1}}}\fm@tmp
  \colorlet{cursfosc}{curs!72!black}%
  \def\fm@logo{\fm@logodir/vertical-#1.pdf}%
  \def\fm@marca{\fm@logodir/marca-#1.pdf}}

% ---------- capçalera de la primera pàgina ----------
% #1 matèria (petit), #2 curs, #3 tipus de document i unitat, #4 títol
\newcommand{\capcalera}[4]{%
  \gdef\fm@cursnom{#2}%
  \par\noindent
  \begin{tcolorbox}[colback=curs,colframe=curs,boxrule=0pt,arc=2.2mm,width=3.5cm,height=2.15cm,
      valign=center,left=3mm,right=2mm,top=1mm,bottom=1mm,nobeforeafter,box align=center]
    \color{curstext}{\footnotesize #1}\par\vspace{1pt}{\fontsize{19}{21}\selectfont\bfseries #2\par}%
  \end{tcolorbox}\hspace{5mm}%
  \begin{minipage}[c]{\dimexpr\linewidth-3.5cm-5mm-2.9cm\relax}
    {\color{cursfosc}\small\bfseries\MakeUppercase{#3}}\par\vspace{2pt}
    {\color{navy}\fontsize{20}{23}\selectfont\bfseries #4\par}
  \end{minipage}\hfill
  \begin{minipage}[c]{2.6cm}\raggedleft\includegraphics[width=2.6cm]{\fm@logo}\end{minipage}\par
  \vspace{2pt}}

% ---------- barres de secció del color del curs ----------
\renewcommand{\seccio}[2]{\par\needspace{\fmneed\baselineskip}\bigskip
  \begin{tcolorbox}[colback=curs,colframe=curs,coltext=curstext,arc=2mm,boxrule=0pt,left=3mm,right=3mm,top=1.2mm,bottom=1.2mm]
  {\large\bfseries #1\quad #2}
  \end{tcolorbox}\vspace{-2pt}}

% ---------- capçalera i peu de les pàgines ----------
% primera pàgina: logos institucionals al peu
\fancypagestyle{primera}{\fancyhf{}\renewcommand{\headrulewidth}{0pt}
  \fancyfoot[L]{\raisebox{-0.32\height}{\includegraphics[height=0.78cm]{logo-gva.png}}\hspace{5mm}%
    \raisebox{-0.32\height}{\includegraphics[height=0.78cm]{logo-vermellar.png}}}
  \fancyfoot[R]{\scriptsize\color{muted}\href{https://creativecommons.org/licenses/by-nc-sa/4.0/deed.ca}{\ccbyncsa}\ \ CC BY-NC-SA 4.0 · fentmates.org\qquad\small\thepage}}
% la resta: el símbol a l'esquerra i el curs a la dreta, amb el títol al centre com abans
\fancyhead[L]{\raisebox{-0.3\height}{\includegraphics[height=0.62cm]{\fm@marca}}}
\fancyhead[R]{\small\bfseries\color{cursfosc}\fm@cursnom}
\renewcommand{\headrule}{\color{curs!45}\hrule height \headrulewidth}
\makeatother

%   \fmcurs{1eso}
% i, al començament del document:
%   \capcalera{Matemàtiques}{1r ESO}{Fitxa d'exercicis · Unitat 2}{Divisibilitat}
% Opcional abans de \input: \def\fmneed{14} (espai mínim abans de cada secció).
% =====================================================================
\makeatletter
\def\fm@logodir{../../logo/pdf}
\providecommand{\fmneed}{10}
\def\fm@cursnom{}

% clau del curs, color i color del text damunt del color
\newcommand{\fm@curs}[3]{\@namedef{fm@color@#1}{#2}\@namedef{fm@text@#1}{#3}}
\fm@curs{1eso}{E6B84A}{navy}
\fm@curs{2eso}{3B82F6}{white}
\fm@curs{3eso}{FF7A50}{white}
\fm@curs{4eso-a}{A78BFA}{white}
\fm@curs{4eso-b}{6D28D9}{white}
\fm@curs{1batx-mi}{9B1C1C}{white}
\fm@curs{1batx-ccss}{EF6B6B}{white}
\fm@curs{2batx-mii}{1E6B47}{white}
\fm@curs{2batx-ccss}{2BB3A3}{white}

\newcommand{\fmcurs}[1]{%
  \@ifundefined{fm@color@#1}{\PackageError{fentmates}{Curs desconegut: #1}{}}{}%
  \edef\fm@tmp{\noexpand\definecolor{curs}{HTML}{\@nameuse{fm@color@#1}}%
               \noexpand\colorlet{curstext}{\@nameuse{fm@text@#1}}}\fm@tmp
  \colorlet{cursfosc}{curs!72!black}%
  \def\fm@logo{\fm@logodir/vertical-#1.pdf}%
  \def\fm@marca{\fm@logodir/marca-#1.pdf}}

% ---------- capçalera de la primera pàgina ----------
% #1 matèria (petit), #2 curs, #3 tipus de document i unitat, #4 títol
\newcommand{\capcalera}[4]{%
  \gdef\fm@cursnom{#2}%
  \par\noindent
  \begin{tcolorbox}[colback=curs,colframe=curs,boxrule=0pt,arc=2.2mm,width=3.5cm,height=2.15cm,
      valign=center,left=3mm,right=2mm,top=1mm,bottom=1mm,nobeforeafter,box align=center]
    \color{curstext}{\footnotesize #1}\par\vspace{1pt}{\fontsize{19}{21}\selectfont\bfseries #2\par}%
  \end{tcolorbox}\hspace{5mm}%
  \begin{minipage}[c]{\dimexpr\linewidth-3.5cm-5mm-2.9cm\relax}
    {\color{cursfosc}\small\bfseries\MakeUppercase{#3}}\par\vspace{2pt}
    {\color{navy}\fontsize{20}{23}\selectfont\bfseries #4\par}
  \end{minipage}\hfill
  \begin{minipage}[c]{2.6cm}\raggedleft\includegraphics[width=2.6cm]{\fm@logo}\end{minipage}\par
  \vspace{2pt}}

% ---------- barres de secció del color del curs ----------
\renewcommand{\seccio}[2]{\par\needspace{\fmneed\baselineskip}\bigskip
  \begin{tcolorbox}[colback=curs,colframe=curs,coltext=curstext,arc=2mm,boxrule=0pt,left=3mm,right=3mm,top=1.2mm,bottom=1.2mm]
  {\large\bfseries #1\quad #2}
  \end{tcolorbox}\vspace{-2pt}}

% ---------- capçalera i peu de les pàgines ----------
% primera pàgina: logos institucionals al peu
\fancypagestyle{primera}{\fancyhf{}\renewcommand{\headrulewidth}{0pt}
  \fancyfoot[L]{\raisebox{-0.32\height}{\includegraphics[height=0.78cm]{logo-gva.png}}\hspace{5mm}%
    \raisebox{-0.32\height}{\includegraphics[height=0.78cm]{logo-vermellar.png}}}
  \fancyfoot[R]{\scriptsize\color{muted}\href{https://creativecommons.org/licenses/by-nc-sa/4.0/deed.ca}{\ccbyncsa}\ \ CC BY-NC-SA 4.0 · fentmates.org\qquad\small\thepage}}
% la resta: el símbol a l'esquerra i el curs a la dreta, amb el títol al centre com abans
\fancyhead[L]{\raisebox{-0.3\height}{\includegraphics[height=0.62cm]{\fm@marca}}}
\fancyhead[R]{\small\bfseries\color{cursfosc}\fm@cursnom}
\renewcommand{\headrule}{\color{curs!45}\hrule height \headrulewidth}
\makeatother

\fmcurs{3eso}

\begin{document}
\thispagestyle{primera}

\capcalera{Matemàtiques}{3r ESO}{Fitxa d'exercicis · Unitat 10}{Proporcionalitat. Teorema de Tales}
\vspace{10pt}
Nom i cognoms: \rule{8.4cm}{0.4pt}\hfill Grup: \rule{1.6cm}{0.4pt}\hfill Data: \rule{2.2cm}{0.4pt}
\vspace{8pt}

\begin{tcolorbox}[colback=soft,colframe=soft,arc=2mm,left=3mm,right=3mm,top=1.5mm,bottom=1.5mm]
\small Exercicis graduats per nivells de dificultat:\quad \estrelles{1}\ bàsic \qquad \estrelles{2}\ mitjà \qquad \estrelles{3}\ avançat.\\[2pt] Resol-los al quadern mostrant tots els passos i fes un dibuix de cada situació. Escriu la proporció abans de calcular i, si un resultat no és exacte, arrodoneix-lo a les centèsimes.
\end{tcolorbox}

% =====================================================
\seccio{1}{Raons i proporcions entre segments}

\ex{1} Calcula la raó entre els segments (si no és exacta, arrodoneix a les centèsimes):\begin{tasks}(4)
  \task 12 cm i 8 cm
  \task 5 cm i 20 cm
  \task 0,6 m i 45 cm
  \task 2 dm i 25 cm
\end{tasks}\recorda{La \textbf{raó} de dos segments és el quocient de les seues longituds, mesurades amb la \textbf{mateixa unitat}. Quatre segments són \textbf{proporcionals} si $\frac ab=\frac cd$, és a dir, si $a\cdot d=b\cdot c$ (producte dels extrems = producte dels mitjans).}

\ex{1} Calcula $x$ en cada proporció:\begin{tasks}(3)
  \task $\dfrac35=\dfrac{x}{20}$
  \task $\dfrac8x=\dfrac27$
  \task $\dfrac x4=\dfrac{15}{6}$
  \task $\dfrac{2{,}5}{x}=\dfrac58$
  \task $\dfrac69=\dfrac{10}{x}$
  \task $\dfrac{x}{12}=\dfrac{0{,}75}{3}$
\end{tasks}

\ex{1} Són proporcionals, en aquest ordre, els segments següents?\begin{tasks}(2)
  \task 2, 3, 4 i 6 cm
  \task 3, 5, 6 i 9 cm
  \task 1,5, 2, 4,5 i 6 cm
  \task 4, 7, 8 i 15 cm
\end{tasks}

\ex{2} Calcula el quart proporcional dels segments: \ a) 4, 6 i 10 cm; \ b) 2,5, 3 i 7 cm.

\ex{2} Dos segments estan en raó $\frac35$. Si el menut fa 12 cm, quant fa el gran?

\ex{2} Dos segments sumen 56 cm i estan en raó $\frac34$. Quant fa cadascun?

\ex{3} El \textbf{mitjà proporcional} de $a$ i $b$ és el segment $x$ tal que $\frac ax=\frac xb$. Calcula el mitjà proporcional de: \ a) 4 i 25; \ b) 2 i 8; \ c) 3 i 6 (arrodoneix a les centèsimes).


% =====================================================
\seccio{2}{El teorema de Tales}

\ex{1} Les rectes \textcolor{navy}{blau fosc} són paral·leles. Calcula $x$.\par\smallskip\begin{center}\begin{tikzpicture}[x=0.55cm,y=0.55cm,baseline=(current bounding box.center),line join=round]
\draw[draw=navy,thick] (-1.200,0.000) -- (5.503,0.000);
\draw[draw=navy,thick] (-1.200,2.000) -- (5.503,2.000);
\draw[draw=navy,thick] (-1.200,5.000) -- (5.503,5.000);
\draw[thick] (-0.490,-0.700) -- (3.991,5.700);
\draw[thick] (3.502,-0.700) -- (4.401,5.700);
\draw[draw=accent,line width=1.6pt] (3.881,2.000) -- (4.303,5.000);
\fill (0.000,0.000) circle (1.5pt);
\fill (1.400,2.000) circle (1.5pt);
\fill (3.501,5.000) circle (1.5pt);
\fill (3.600,0.000) circle (1.5pt);
\fill (3.881,2.000) circle (1.5pt);
\fill (4.303,5.000) circle (1.5pt);
\node[font=\small,inner sep=2pt,above left] at (0.000,0.000) {$A$};
\node[font=\small,inner sep=2pt,above left] at (1.400,2.000) {$B$};
\node[font=\small,inner sep=2pt,above left] at (3.501,5.000) {$C$};
\node[font=\small,inner sep=2pt,above right] at (3.600,0.000) {$A'$};
\node[font=\small,inner sep=2pt,above right] at (3.881,2.000) {$B'$};
\node[font=\small,inner sep=2pt,above right] at (4.303,5.000) {$C'$};
\node[font=\small,inner sep=2pt,left] at (0.700,1.000) {$4$};
\node[font=\small,inner sep=2pt,left] at (2.451,3.500) {$6$};
\node[font=\small,inner sep=2pt,right] at (3.741,1.000) {$6$};
\node[font=\small,inner sep=2pt,right,text=accent] at (4.092,3.500) {$x$};
\end{tikzpicture}\end{center}\recorda{\textbf{Teorema de Tales:} si unes rectes paral·leles tallen dues secants, els segments que determinen sobre una són proporcionals als de l'altra: $\frac{\overline{AB}}{\overline{BC}}=\frac{\overline{A'B'}}{\overline{B'C'}}$, o bé $\frac{\overline{AB}}{\overline{A'B'}}=\frac{\overline{BC}}{\overline{B'C'}}=\frac{\overline{AC}}{\overline{A'C'}}$.}

\ex{1} Les rectes $r$, $s$ i $t$ són paral·leles i tallen les secants en $A,\,B,\,C$ i $A',\,B',\,C'$. Calcula el segment que falta:\begin{tasks}(1)
  \task $\overline{AB}=3$, $\overline{BC}=5$, $\overline{A'B'}=4{,}8$. Calcula $\overline{B'C'}$.
  \task $\overline{AB}=6$, $\overline{A'B'}=4$, $\overline{B'C'}=5$. Calcula $\overline{BC}$.
  \task $\overline{AB}=4$, $\overline{AC}=10$, $\overline{A'C'}=15$. Calcula $\overline{A'B'}$.
\end{tasks}

\ex{2} Les rectes \textcolor{navy}{blau fosc} són paral·leles. Calcula $x$.\par\smallskip\begin{center}\begin{tikzpicture}[x=0.55cm,y=0.55cm,baseline=(current bounding box.center),line join=round]
\draw[draw=navy,thick] (-1.200,0.000) -- (5.475,0.000);
\draw[draw=navy,thick] (-1.200,1.600) -- (5.475,1.600);
\draw[draw=navy,thick] (-1.200,4.800) -- (5.475,4.800);
\draw[thick] (-0.490,-0.700) -- (3.851,5.500);
\draw[thick] (3.502,-0.700) -- (4.373,5.500);
\draw[draw=accent,line width=1.6pt] (0.000,0.000) -- (1.120,1.600);
\fill (0.000,0.000) circle (1.5pt);
\fill (1.120,1.600) circle (1.5pt);
\fill (3.361,4.800) circle (1.5pt);
\fill (3.600,0.000) circle (1.5pt);
\fill (3.825,1.600) circle (1.5pt);
\fill (4.275,4.800) circle (1.5pt);
\node[font=\small,inner sep=2pt,above left] at (0.000,0.000) {$A$};
\node[font=\small,inner sep=2pt,above left] at (1.120,1.600) {$B$};
\node[font=\small,inner sep=2pt,above left] at (3.361,4.800) {$C$};
\node[font=\small,inner sep=2pt,above right] at (3.600,0.000) {$A'$};
\node[font=\small,inner sep=2pt,above right] at (3.825,1.600) {$B'$};
\node[font=\small,inner sep=2pt,above right] at (4.275,4.800) {$C'$};
\node[font=\small,inner sep=2pt,left,text=accent] at (0.560,0.800) {$x$};
\node[font=\small,inner sep=2pt,left] at (2.241,3.200) {$6$};
\node[font=\small,inner sep=2pt,right] at (3.712,0.800) {$4$};
\node[font=\small,inner sep=2pt,right] at (4.050,3.200) {$8$};
\end{tikzpicture}\end{center}

\ex{2} Les rectes que tallen dues secants hi determinen aquests segments. Són paral·leles?\begin{tasks}(3)
  \task 2 i 3; \ 4 i 6
  \task 3 i 4; \ 4 i 5
  \task 1,5 i 2,5; \ 2,4 i 4
\end{tasks}

\ex{2} Tres parcel·les estan entre dos carrers i les seues partions són paral·leles. Al carrer de la Mar fan 15 m, 20 m i 25 m de façana. Al carrer Ample, les tres juntes fan 72 m. Quant fa la façana de cada parcel·la al carrer Ample?

\ex{3} Unes rectes paral·leles determinen $\overline{AB}=x$ i $\overline{BC}=x+4$ sobre una secant, i $\overline{A'B'}=3$ i $\overline{B'C'}=5$ sobre l'altra. Calcula $x$.

\ex{3} Unes rectes paral·leles determinen sobre una secant $\overline{AB}=5$ cm i $\overline{BC}=8$ cm. Sobre una altra secant, $\overline{A'C'}=19{,}5$ cm. Calcula $\overline{A'B'}$ i $\overline{B'C'}$.


% =====================================================
\seccio{3}{Dividir un segment en parts iguals o proporcionals}

\ex{1} Dibuixa un segment de 12 cm i divideix-lo en 5 parts iguals amb regle i compàs. Quant fa cada part?\recorda{Des d'un extrem, traça una semirecta auxiliar i marca-hi amb el compàs tants segments iguals com parts vulgues. Uneix l'última marca amb l'altre extrem i traça paral·leles per les altres marques. Per a parts proporcionals a $2,\,3,\,4$, marca $2+3+4=9$ segments i traça les paral·leles per les marques 2 i 5.}

\ex{1} Calcula les parts en què queda dividit:\begin{tasks}(1)
  \task un segment de 30 cm, en parts proporcionals a 2 i 3
  \task un segment de 24 cm, en parts proporcionals a 1, 2 i 3
  \task un segment de 45 cm, en parts proporcionals a 2, 3 i 4
\end{tasks}

\ex{2} Dibuixa un segment de 10 cm i divideix-lo, amb regle i compàs, en dues parts proporcionals a 2 i 3. Explica els passos i comprova el resultat calculant.

\ex{2} Un punt $P$ divideix un segment $AB$ de 35 cm de manera que $\frac{\overline{AP}}{\overline{PB}}=\frac25$. Quant fan $\overline{AP}$ i $\overline{PB}$?

\ex{2} Tallem una corda de 3,6 m en tres trossos proporcionals a 1, 2 i 5. Quant fa cada tros?

\ex{3} Un punt $P$ divideix un segment $AB$ en dues parts que estan en raó $\frac37$, i la part llarga fa 12 cm més que la curta. Quant fan les dues parts i el segment sencer?


% =====================================================
\seccio{4}{Triangles en posició de Tales}

\ex{1} En la figura, $B'C'\parallel BC$. Calcula $x$.\par\smallskip\begin{center}\begin{tikzpicture}[x=0.5cm,y=0.5cm,baseline=(current bounding box.center),line join=round]
\draw[fill=teal!12,draw=teal,thick] (0.000,0.000) -- (-2.817,-5.298) -- (2.644,-5.938) -- cycle;
\draw[fill=accent!22,draw=accent,thick] (0.000,0.000) -- (-1.127,-2.119) -- (1.058,-2.375) -- cycle;
\draw[draw=accent,line width=1.6pt] (-1.127,-2.119) -- (1.058,-2.375);
\draw[draw=teal,line width=1.6pt] (-2.817,-5.298) -- (2.644,-5.938);
\fill (0.000,0.000) circle (1.5pt);
\fill (-2.817,-5.298) circle (1.5pt);
\fill (2.644,-5.938) circle (1.5pt);
\fill (-1.127,-2.119) circle (1.5pt);
\fill (1.058,-2.375) circle (1.5pt);
\node[font=\small,inner sep=2pt,above] at (0.000,0.000) {$A$};
\node[font=\small,inner sep=2pt,left] at (-2.817,-5.298) {$B$};
\node[font=\small,inner sep=2pt,right] at (2.644,-5.938) {$C$};
\node[font=\small,inner sep=2pt,left] at (-1.127,-2.119) {$B'$};
\node[font=\small,inner sep=2pt,right] at (1.058,-2.375) {$C'$};
\node[font=\small,inner sep=2pt,left] at (-0.563,-1.060) {$4$};
\node[font=\small,inner sep=2pt,left] at (-1.972,-3.708) {$6$};
\node[font=\small,inner sep=2pt,below] at (-0.087,-5.618) {$15$};
\node[font=\small,inner sep=2pt,above,text=accent] at (-0.035,-2.247) {$x$};
\end{tikzpicture}\end{center}\recorda{Si $B'C'\parallel BC$, els triangles $AB'C'$ i $ABC$ estan en \textbf{posició de Tales}: $\frac{\overline{AB'}}{\overline{AB}}=\frac{\overline{AC'}}{\overline{AC}}=\frac{\overline{B'C'}}{\overline{BC}}$. Per al costat paral·lel, compara sempre \textbf{costats sencers} des del vèrtex $A$.}

\ex{1} En un triangle $ABC$, $B'$ és un punt de $AB$ i $C'$ un punt de $AC$, amb $B'C'\parallel BC$. Calcula:\begin{tasks}(1)
  \task $\overline{B'C'}$, si $\overline{AB'}=3$, $\overline{B'B}=6$ i $\overline{BC}=12$.
  \task $\overline{BC}$, si $\overline{AB'}=5$, $\overline{AB}=8$ i $\overline{B'C'}=7{,}5$.
  \task $\overline{BC}$, si $\overline{AC'}=4$, $\overline{C'C}=2$ i $\overline{B'C'}=6$.
\end{tasks}

\ex{2} En la figura, $B'C'\parallel BC$. Calcula $x$ i $y$.\par\smallskip\begin{center}\begin{tikzpicture}[x=0.5cm,y=0.5cm,baseline=(current bounding box.center),line join=round]
\draw[fill=teal!12,draw=teal,thick] (0.000,0.000) -- (-2.817,-5.298) -- (2.644,-5.938) -- cycle;
\draw[fill=accent!22,draw=accent,thick] (0.000,0.000) -- (-1.690,-3.179) -- (1.586,-3.563) -- cycle;
\draw[draw=accent,line width=1.6pt] (-1.690,-3.179) -- (1.586,-3.563);
\draw[draw=teal,line width=1.6pt] (-2.817,-5.298) -- (2.644,-5.938);
\fill (0.000,0.000) circle (1.5pt);
\fill (-2.817,-5.298) circle (1.5pt);
\fill (2.644,-5.938) circle (1.5pt);
\fill (-1.690,-3.179) circle (1.5pt);
\fill (1.586,-3.563) circle (1.5pt);
\node[font=\small,inner sep=2pt,above] at (0.000,0.000) {$A$};
\node[font=\small,inner sep=2pt,left] at (-2.817,-5.298) {$B$};
\node[font=\small,inner sep=2pt,right] at (2.644,-5.938) {$C$};
\node[font=\small,inner sep=2pt,left] at (-1.690,-3.179) {$B'$};
\node[font=\small,inner sep=2pt,right] at (1.586,-3.563) {$C'$};
\node[font=\small,inner sep=2pt,left] at (-0.845,-1.589) {$6$};
\node[font=\small,inner sep=2pt,left] at (-2.253,-4.238) {$4$};
\node[font=\small,inner sep=2pt,right] at (0.793,-1.781) {$9$};
\node[font=\small,inner sep=2pt,right] at (2.115,-4.750) {$y$};
\node[font=\small,inner sep=2pt,above] at (-0.052,-3.371) {$7{,}2$};
\node[font=\small,inner sep=2pt,below] at (-0.087,-5.618) {$x$};
\end{tikzpicture}\end{center}

\ex{2} Un triangle té 18 cm de base i 12 cm d'altura. Tracem una paral·lela a la base a 4 cm del vèrtex oposat. Quant mesura el segment de la paral·lela que queda dins del triangle?

\ex{2} En la figura, $AB\parallel DE$. Calcula $x=\overline{DE}$ i $y=\overline{CD}$.\par\smallskip\begin{center}\begin{tikzpicture}[x=0.42cm,y=0.42cm,baseline=(current bounding box.center),line join=round]
\draw[fill=teal!12,draw=teal,thick] (0.000,0.000) -- (6.000,0.000) -- (2.400,3.200) -- cycle;
\draw[fill=accent!22,draw=accent,thick] (2.400,3.200) -- (6.000,8.000) -- (-3.000,8.000) -- cycle;
\node[font=\small,below left] at (0.000,0.000) {$A$};
\node[font=\small,below right] at (6.000,0.000) {$B$};
\node[font=\small,left=4pt] at (2.400,3.200) {$C$};
\node[font=\small,above left] at (-3.000,8.000) {$D$};
\node[font=\small,above right] at (6.000,8.000) {$E$};
\node[font=\small,below] at (3.000,0.000) {$8$};
\node[font=\small,left] at (1.200,1.600) {$6$};
\node[font=\small,right] at (4.200,1.600) {$5$};
\node[font=\small,above] at (1.500,8.000) {$x$};
\node[font=\small,right] at (4.200,5.600) {$9$};
\node[font=\small,left] at (-0.300,5.600) {$y$};
\end{tikzpicture}\end{center}

\ex{3} Per a mesurar l'amplària d'un riu, l'Àlex clava una estaca $B$ a la vora, davant d'un arbre $A$ de l'altra riba. Camina 15 m perpendicularment a $AB$ fins a $C$, en clava una altra, continua 5 m fins a $D$ i gira 90° allunyant-se del riu fins que veu $C$ i $A$ alineats, en el punt $E$, que està a 4 m de $D$. Quina amplària té el riu?

\ex{3} En un triangle $ABC$, $B'C'\parallel BC$, $\overline{B'B}=6$, $\overline{B'C'}=4$ i $\overline{BC}=10$. Calcula $x=\overline{AB'}$.


% =====================================================
\seccio{5}{Figures semblants. Raó de semblança}

\ex{1} Són semblants aquestes parelles de rectangles? Si ho són, calcula la raó de semblança.\begin{tasks}(2)
  \task $4\times6$ \ i \ $6\times9$
  \task $2\times5$ \ i \ $4\times8$
  \task $3\times4$ \ i \ $6\times8$
  \task un quadrat de 3 cm i un de 7 cm
\end{tasks}\recorda{Dues figures són \textbf{semblants} si tenen els angles iguals i els costats homòlegs proporcionals. La \textbf{raó de semblança} és $k=\frac{a'}{a}$. Els perímetres estan en raó $k$ i les àrees en raó $k^2$.}

\ex{1} Un triangle té costats de 4, 6 i 7 cm. Calcula els costats d'un triangle semblant amb raó: \ a) $k=2{,}5$; \ b) $k=0{,}4$.

\ex{2} Una foto de 9 cm × 13 cm s'amplia fins que el costat menut fa 27 cm. Quant fa el costat gran? Quantes vegades més gran és l'àrea?

\ex{2} Dos pentàgons són semblants amb raó $k=1{,}5$. El menut té 20 cm de perímetre i 27,5 cm² d'àrea. Calcula el perímetre i l'àrea del gran.

\ex{2} Dues figures semblants tenen àrees de 18 cm² i 72 cm². Quina és la raó de semblança? Si el perímetre de la menuda és 16 cm, quin és el de la gran?

\ex{3} Tenim un rectangle de 8 cm × 12 cm i en volem dibuixar un de semblant amb una àrea de 600 cm². Quines dimensions ha de tenir?

\ex{3} Un triangle rectangle té catets de 3 cm i 4 cm. Un triangle semblant té 54 cm² d'àrea. Quant fan els seus costats?


% =====================================================
\seccio{6}{Criteris de semblança de triangles}

\ex{1} Són semblants els triangles que tenen aquests costats? Si ho són, dona la raó de semblança.\begin{tasks}(2)
  \task 2, 3, 4 \ i \ 6, 9, 12
  \task 3, 4, 6 \ i \ 6, 8, 10
  \task 5, 5, 8 \ i \ 7,5, 7,5, 12
  \task 4, 6, 9 \ i \ 6, 9, 12
\end{tasks}\recorda{Dos triangles són semblants si tenen: \textbf{1r)} dos angles iguals; \ o \textbf{2n)} els tres costats proporcionals; \ o \textbf{3r)} dos costats proporcionals i l'angle comprés igual. Per a comparar costats, ordena'ls de menor a major.}

\ex{1} Són semblants aquests triangles? Justifica-ho.\begin{tasks}(1)
  \task Un té angles de 35° i 75°, i l'altre, de 75° i 70°.
  \task Un té angles de 50° i 60°, i l'altre, de 60° i 80°.
  \task Dos triangles rectangles, l'un amb un angle de 30° i l'altre amb un angle de 60°.
\end{tasks}

\ex{2} Els triangles $ABC$ i $A'B'C'$ tenen $\hat A=\hat A'=50^\circ$, $\overline{AB}=4$, $\overline{AC}=6$, $\overline{A'B'}=6$ i $\overline{A'C'}=9$. Són semblants? Si $\overline{BC}=5$, quant fa $\overline{B'C'}$?

\ex{2} Un triangle té costats de 6, 8 i 10 cm. El costat menut d'un triangle semblant fa 9 cm. Calcula els altres dos costats i la raó entre les àrees.

\ex{2} Són semblants dos triangles isòsceles que tenen l'angle desigual de 40°? I dos triangles isòsceles qualssevol?

\ex{2} Els catets d'un triangle rectangle fan 6 cm i 8 cm. Calcula la hipotenusa i, fent servir triangles semblants, l'altura sobre la hipotenusa.

\ex{3} Els catets d'un triangle rectangle $ABC$ (angle recte en $A$) fan $\overline{AB}=9$ cm i $\overline{AC}=12$ cm. L'altura $AH$ divideix la hipotenusa en dos segments. Calcula $\overline{BC}$, $\overline{BH}$, $\overline{HC}$ i $\overline{AH}$.

\ex{3} En un trapezi $ABCD$ les bases fan $\overline{AB}=12$ cm i $\overline{CD}=8$ cm, i les diagonals es tallen en $O$. Si $\overline{AO}=9$ cm, quant fa $\overline{OC}$?


% =====================================================
\seccio{7}{Aplicacions: altures i distàncies}

\ex{1} Un pal d'1 m fa una ombra d'1,25 m. Al mateix moment, un arbre fa una ombra de 10 m. Quina altura té l'arbre?\par\smallskip\begin{center}\begin{tikzpicture}[x=0.32cm,y=0.32cm,baseline=(current bounding box.center),line join=round]
\draw[muted] (-0.600,0.000) -- (17.000,0.000);
\draw[fill=accent!22,draw=accent,thick] (0.000,0.000) -- (0.000,2.000) -- (2.500,0.000) -- cycle;
\draw[thin] (0.688,0.000) -- (0.688,0.688) -- (0.000,0.688);
\draw[fill=teal!12,draw=teal,thick] (6.200,0.000) -- (6.200,8.000) -- (16.200,0.000) -- cycle;
\draw[thin] (6.888,0.000) -- (6.888,0.688) -- (6.200,0.688);
\draw[line width=2pt,draw=muted] (0.000,0.000) -- (0.000,2.000);
\draw[line width=2pt,draw=muted] (6.200,0.000) -- (6.200,8.000);
\draw[dashed,draw=muted] (0.000,2.000) -- (2.500,0.000);
\draw[dashed,draw=muted] (6.200,8.000) -- (16.200,0.000);
\node[font=\small,inner sep=2pt,left] at (0.000,1.000) {$1$\,m};
\node[font=\small,inner sep=2pt,below] at (1.250,0.000) {$1{,}25$\,m};
\node[font=\small,inner sep=2pt,left,text=accent] at (6.200,4.000) {$H$};
\node[font=\small,inner sep=2pt,below] at (11.200,0.000) {$10$\,m};
\end{tikzpicture}\end{center}\recorda{Busca dos \textbf{triangles semblants} (rajos de sol paral·lels, un mirall, una visual que passa per dos punts…) i escriu la proporció entre costats homòlegs: $\frac{\text{altura}}{\text{ombra}}=\frac{\text{altura}'}{\text{ombra}'}$.}

\ex{1} Un edifici de 24 m fa una ombra de 18 m. Quina ombra fa, al mateix moment, una persona d'1,6 m?

\ex{2} La Marta té els ulls a 1,5 m de terra. Posa un mirall $M$ a terra a 2 m d'ella i veu reflectida la punta d'una torre que està a 18 m del mirall. Quina altura té la torre?\par\smallskip\begin{center}\begin{tikzpicture}[x=0.26cm,y=0.26cm,baseline=(current bounding box.center),line join=round]
\draw[muted] (-0.500,0.000) -- (21.000,0.000);
\draw[line width=2pt,draw=muted] (0.000,0.000) -- (0.000,1.500);
\draw[line width=2pt,draw=muted] (20.000,0.000) -- (20.000,13.500);
\draw[dashed,draw=muted] (0.000,1.500) -- (2.000,0.000);
\draw[dashed,draw=muted] (2.000,0.000) -- (20.000,13.500);
\fill (2.000,0.000) circle (1.5pt);
\node[font=\small,left] at (0.000,0.750) {$1{,}5$\,m};
\node[font=\small,below] at (1.000,-1.600) {$2$\,m};
\node[font=\small,below] at (11.000,-1.600) {$18$\,m};
\node[font=\small,right] at (20.000,6.750) {$H$};
\node[font=\small,below] at (2.000,0.000) {$M$};
\end{tikzpicture}\end{center}

\ex{2} Tales va mesurar la gran piràmide amb les ombres. Si en un moment un bastó de 2 m fa una ombra de 2 m i l'ombra de la piràmide, comptada des del centre de la base, fa 146 m, quina altura té la piràmide?

\ex{2} Una estaca d'1,5 m està clavada a 2 m d'un punt $P$ de terra. Des de $P$, la punta de l'estaca i la punta d'una torre que està a 30 m de $P$ queden alineades. Quina altura té la torre?

\ex{3} Una persona d'1,6 m està a 3 m d'un fanal de 5 m d'altura. Quant fa la seua ombra?\par\smallskip\begin{center}\begin{tikzpicture}[x=0.55cm,y=0.55cm,baseline=(current bounding box.center),line join=round]
\draw[muted] (-0.500,0.000) -- (5.200,0.000);
\draw[line width=2pt,draw=muted] (0.000,0.000) -- (0.000,5.000);
\draw[line width=2pt,draw=muted] (3.000,0.000) -- (3.000,1.600);
\draw[dashed,draw=muted] (0.000,5.000) -- (4.412,0.000);
\fill (0.000,5.000) circle (1.5pt);
\node[font=\small,left] at (0.000,2.500) {$5$\,m};
\node[font=\small,left] at (3.000,0.800) {$1{,}6$\,m};
\node[font=\small,below] at (1.500,0.000) {$3$\,m};
\node[font=\small,below] at (3.706,0.000) {$s$};
\end{tikzpicture}\end{center}

\ex{3} Per a mesurar l'amplària $\overline{AB}$ d'un riu, des de $B$ (davant d'un arbre $A$ de l'altra riba) caminem 20 m perpendicularment a $AB$ fins a $C$ i 8 m més fins a $D$. Des de $D$ ens allunyem del riu, perpendicularment, fins al punt $E$ des d'on veiem $C$ i $A$ alineats. Si $\overline{DE}=6$ m, quina amplària té el riu?

\ex{3} Una persona amb els ulls a 1,5 m de terra està a 3 m d'una tàpia de 3 m. Veu la part de dalt de la tàpia alineada amb la part de dalt d'un edifici que està 7 m darrere de la tàpia. Quina altura té l'edifici?\par\smallskip\begin{center}\begin{tikzpicture}[x=0.45cm,y=0.45cm,baseline=(current bounding box.center),line join=round]
\draw[muted] (-0.500,0.000) -- (10.800,0.000);
\draw[line width=2pt,draw=muted] (0.000,0.000) -- (0.000,1.500);
\draw[line width=2pt,draw=muted] (3.000,0.000) -- (3.000,3.000);
\draw[line width=2pt,draw=muted] (10.000,0.000) -- (10.000,6.500);
\draw[dashed,draw=muted] (0.000,1.500) -- (10.000,6.500);
\draw[dashed,draw=muted] (0.000,1.500) -- (10.000,1.500);
\node[font=\small,left] at (0.000,0.750) {$1{,}5$};
\node[font=\small,right] at (3.000,1.500) {$3$\,m};
\node[font=\small,below] at (1.500,0.000) {$3$\,m};
\node[font=\small,below] at (6.500,0.000) {$7$\,m};
\node[font=\small,right] at (10.000,4.000) {$H$};
\end{tikzpicture}\end{center}


% =====================================================
\seccio{8}{Escales. Àrees i volums de figures semblants}

\ex{1} Calcula la distància real:\begin{tasks}(3)
  \task 8 cm en un plànol 1:50
  \task 6 cm en un mapa 1:25\,000
  \task 3,5 cm en un mapa 1:200\,000
\end{tasks}\recorda{Escala $1:n$: 1 cm del plànol són $n$ cm reals. \ Real $=$ plànol $\times\,n$; \ plànol $=$ real $:\,n$. Si la raó de semblança és $k$, les àrees es multipliquen per $k^2$ i els volums per $k^3$.}

\ex{1} Quina longitud tindrà en el plànol? \ a) Una paret de 12 m, en un plànol 1:100. \ b) Una carretera de 45 km, en un mapa 1:1\,500\,000.

\ex{2} Calcula l'escala: \ a) 5 cm del mapa representen 2 km; \ b) 3 cm del plànol representen 150 m.

\ex{2} En el plànol d'un pis a escala 1:200, el menjador fa 3 cm × 2,5 cm. Calcula'n les dimensions reals i la superfície.

\ex{2} L'escala gràfica d'un mapa indica que 2 cm són 500 m. Quina és l'escala numèrica? Quina distància real hi ha entre dos punts que estan a 7 cm en el mapa?

\ex{2} Una maqueta d'un cotxe està feta a escala 1:18. Si el cotxe real fa 4,5 m, quant fa la maqueta? El parabrisa de la maqueta té 40 cm². Quina superfície té el real?

\ex{3} Dues botelles semblants fan 15 cm i 30 cm d'altura. La menuda té una capacitat de 0,25 L i una etiqueta de 60 cm². Quina capacitat té la gran? I quina superfície té la seua etiqueta?

\ex{3} En un mapa a escala 1:50\,000, un bosc ocupa 6 cm². Quina superfície real té, en km² i en hectàrees?

\vfill
\begin{center}\color{muted}\small\faCheckCircle[regular]\ \ Ànim! Si et costa algun exercici, repassa els apunts interactius de la unitat.\end{center}

\end{document}
